\section{Preliminaries}
\label{sec:prelim}

\subsection{Notation}

Throughout, $\lambda$ denotes a security parameter (set to $128$ for
NIST PQ Category 3). For a finite set $S$, $x \xleftarrow{\$} S$ denotes
uniform sampling. A function $\epsilon(\lambda)$ is \emph{negligible}
if for every polynomial $p$, $\epsilon(\lambda) < 1/p(\lambda)$ for
all sufficiently large $\lambda$; we write $\mathsf{negl}(\lambda)$.
The notation $H_X$ denotes a hash function modeled as a random oracle
in domain-separated context $X$. The string $\mathsf{ctx}_X$ denotes
the FIPS 204 \S5 context byte string for purpose $X$, fixed in
Table~\ref{tab:domsep}.

\subsection{Validator set}

A Lux chain at consensus epoch $e$ has a validator set
$\mathcal{V}_e = \{v_1, \ldots, v_n\}$ with stake distribution
$\mathrm{stk}: \mathcal{V}_e \to \mathbb{N}$ and total stake
$S_e = \sum_i \mathrm{stk}(v_i)$. Each validator $v_i$ holds:

\begin{itemize}[leftmargin=2em,nosep]
\item a BLS12-381 keypair $(\mathit{sk}_i^{\mathrm{BLS}},\,
      \mathit{pk}_i^{\mathrm{BLS}})$ with proof of possession;
\item a Pulsar threshold-secret share $\mathit{sk}_i^{\mathrm{Puls}}$
      over the joint Pulsar group public key
      $\mathit{pk}^{\mathrm{Puls}}$, produced by the Pulsar DKG
      \cite{pulsar};
\item on the Aurora and Polaris profiles, a Corona share
      $\mathit{sk}_i^{\mathrm{Cor}}$ over a joint
      $\mathit{pk}^{\mathrm{Cor}}$ \cite{corona};
\item on the Polaris profile, a Magnetar share
      $\mathit{sk}_i^{\mathrm{Mag}}$ over a joint
      $\mathit{pk}^{\mathrm{Mag}}$ \cite{magnetar};
\item a per-validator ML-DSA-65 keypair $(\mathit{sk}_i^{\mathrm{ML}},
      \mathit{pk}_i^{\mathrm{ML}})$ used for the Z-Chain rollup leg.
\end{itemize}

The Byzantine threshold is $f < n/3$ in count and $f < S_e/3$ in
stake. The quorum threshold $\tau = \lfloor 2n/3 \rfloor + 1$ in count
(or $\lceil 2 S_e / 3 \rceil$ in stake) is the minimum signer set
required to assemble a valid leg. The validator set is committed to a
Merkle root $\mathsf{validator\_root}_e$ at epoch boundaries.

\subsection{Round digest}

Each consensus round produces a \emph{round digest} $\mu$ that all legs
sign or prove against. The round digest is defined as

\begin{equation}
\mu \;=\; H_{\mathsf{round}}\!\bigl(\mathsf{chain\_id}\,\|\, e\,\|\, r\,\|\,
\mathsf{validator\_root}_e\,\|\, \mathsf{block\_hash}\bigr),
\end{equation}

where $\mathsf{chain\_id}$ is the 32-byte unique chain identifier, $e$
is the epoch, $r$ is the round number, and $\mathsf{block\_hash}$ is
the structural hash of the block being finalized. The digest is a
32-byte SHA3-256 / Keccak-256 output. This binding makes cross-chain
and cross-epoch replay structurally impossible (Theorem~\ref{thm:replay}).

\subsection{BLS12-381 aggregate signatures}

We use the proof-of-possession variant of BLS aggregation~\cite{ietfbls,
boneh2001short,bonehshorth}: keys are generated as
$\mathit{sk} \xleftarrow{\$} \mathbb{Z}_r$,
$\mathit{pk} = \mathit{sk} \cdot G_2 \in \mathbb{G}_2$, each public
key is published with a proof of possession (a BLS signature on the
public key under itself, in a domain-separated context); signatures
are $\sigma = \mathit{sk} \cdot H_{\mathbb{G}_1}(\mu) \in \mathbb{G}_1$;
aggregation is point addition, $\sigma_{\mathrm{agg}} = \sum_{i \in Q}
\sigma_i$; verification is
$e(G_2, \sigma_{\mathrm{agg}}) = e(\mathit{apk}, H_{\mathbb{G}_1}(\mu))$
with aggregated public key $\mathit{apk} = \sum_{i \in Q} \mathit{pk}_i$.
Aggregated signatures are 48 bytes (compressed $\mathbb{G}_1$);
verification cost is $O(1)$ in $|Q|$ and dominated by two pairings.

\paragraph{Security assumption.}
BLS-PoP is EUF-CMA under the co-CDH assumption on
$(\mathbb{G}_1, \mathbb{G}_2, \mathbb{G}_T)$ in the random oracle
model \cite{boneh2001short,boldyreva2003,bonehshorth}. The co-CDH
assumption is broken in BQP by Shor's algorithm on the discrete
logarithm in $\mathbb{F}_{q^{12}}$, so BLS contributes no PQ
security: it is the \emph{classical} fast-path leg.

\subsection{Pulsar: threshold ML-DSA-65}

Pulsar~\cite{pulsar} is the Lux production fork of the threshold
ML-DSA construction of Celi et al.~\cite{tmldsa2025}, integrated with
the Lux Secret Sharing (LSS) lifecycle~\cite{luxlssmpc}. The signing
output is byte-identical to a single-party FIPS 204
\cite{fips204} ML-DSA-65 signature on the group public key
$\mathit{pk}^{\mathrm{Puls}}$, so verification uses the unmodified
NIST verifier. NIST MPTC \cite{nistir8214c} classifies this as Class
N1 (single-party-compatible threshold signing); the Lux deployment
adds Class N4 (public-key preservation across resharing).

\paragraph{Security assumption.}
Pulsar's unforgeability reduces to the EUF-CMA security of FIPS 204
ML-DSA-65, which in turn reduces to MLWE and MSIS at NIST Category 3
(${\sim}128$-bit quantum security). The v0.1 Pulsar deployment has
an aggregator-as-TCB caveat: a malicious aggregator can refuse to
produce a signature (a liveness fault) but cannot produce a forgery,
because the aggregator never holds enough material to bypass the
threshold; this is documented in the Pulsar deployment runbook
\cite{pulsar}.

\subsection{Corona: threshold M-LWE}

Corona~\cite{corona,ringtail} is a two-round lattice
threshold signature over the polynomial ring $R_q = \mathbb{Z}_q[X]/
(X^N+1)$ with the Lux production parameter set $N = 256$, $q =
2^{48} + 2561$ (NTT-friendly), module dimension $\ell = 8$. It
follows a Schnorr-with-aborts structure: each party samples a nonce
$\mathbf{y}_i$, broadcasts a commitment $\mathbf{w}_i = \mathbf{A}
\mathbf{y}_i$, the challenge is $c = H_{\mathsf{cor}}(\mu \,\|\, \sum_i
\mathbf{w}_i)$, and the response is $\mathbf{z}_i = \mathbf{y}_i + c
\cdot \mathit{sk}_i^{\mathrm{Cor}}$ (with rejection sampling).
Combination is Lagrange interpolation over the ring; verification
checks $\|\mathbf{z}\| \leq B$ and the Fiat--Shamir relation.

\paragraph{Security assumption.}
Corona is SUF-CMA under the Module-LWE and Module-SIS assumptions
in the random oracle model \cite{ringtail}. Corona and ML-DSA both
reduce to Module-LWE; they differ in parameter set and codebase
(Corona at the Ringtail-derived parameters, ML-DSA at the FIPS-204
tuning), not in hardness family. Pairing them gives parameter- and
implementation-level diversity at production sizes
\cite{albrechtestimator,bdgl2016} --- a break specific to one
parameter set need not transfer to the other --- but not
family-disjoint hardness; that is supplied by the hash-based
Magnetar leg.

\subsection{Magnetar: threshold SLH-DSA}

Magnetar~\cite{magnetar} is a threshold construction of FIPS 205
\cite{fips205} SLH-DSA via \emph{reveal-and-aggregate} over the
SLH-DSA scheme seed. The DKG protocol shares the joint master seed
$S$ via byte-wise Shamir secret sharing over $\mathrm{GF}(257)$
\cite{shamir1979}, with $t$-of-$n$ reconstruction; signing
reconstructs $S$ inside an aggregator's memory, calls
\texttt{slhdsa.SignDeterministic} once on the reconstructed seed, and
immediately zeroizes. The signature is byte-identical to a single-party
FIPS 205 SLH-DSA signature on the joint group public key.

\paragraph{Security assumption.}
Magnetar's unforgeability reduces to SLH-DSA's EUF-CMA security,
which in turn reduces to the second-preimage and target-collision
resistance of SHA-3 / SHAKE-192 (or SHAKE-256, depending on
parameter set). The v0.1 reveal-and-aggregate construction has an
aggregator-as-TCB caveat identical in spirit to Pulsar v0.1: the
aggregator briefly holds the reconstructed master seed and must
zeroize. The Tier B status of Magnetar reflects that this caveat
has not yet been audited end-to-end; production deployment of the
Polaris profile is gated on Tier A audit of Magnetar
\cite{magnetar}.

\subsection{Z-Chain Groth16 rollup of ML-DSA-65}

The Z-Chain VM~\cite{chainsrepo} runs a Groth16~\cite{groth2016}
prover for the relation

\begin{equation}
\mathcal{R}_{\mathrm{ZK}} \;=\; \Bigl\{\, (x;\, w)
\;:\; x = (\mu, h_x),\;
w = \bigl(\{\mathit{pk}_i^{\mathrm{ML}}\}_i,
\{\sigma_i^{\mathrm{ML}}\}_i\bigr)\,\Bigr\},
\end{equation}

with predicate

\begin{equation*}
\Phi(x, w) \;\iff\;
h_x = H_{\mathsf{Poseidon}}\!\bigl(\{\mathit{pk}_i^{\mathrm{ML}}\}_i\,\|\,\mu\bigr)
\;\;\wedge\;\;
\sum_i \bigl[\!\bigl[\,\mathsf{MLDSA.Verify}(\mathit{pk}_i^{\mathrm{ML}}, \mu, \sigma_i^{\mathrm{ML}}) = 1\,\bigr]\!\bigr]
\;\geq\; \tau,
\end{equation*}

where $\tau = \lfloor 2n/3 \rfloor + 1$ and ML-DSA verification is
the FIPS 204 algorithm. The Groth16 proof $\pi_{\mathrm{ZK}}$ is
$(A, B, C) \in \mathbb{G}_1 \times \mathbb{G}_2 \times \mathbb{G}_1$,
totalling $48 + 96 + 48 = 192$ bytes compressed.

\paragraph{Security assumption.}
Groth16 is knowledge-sound in the algebraic group model on
BLS12-381 \cite{groth2016}; the verifying key $\mathit{vk}_Z$ is the
output of a powers-of-tau ceremony plus a circuit-specific phase 2
with $\geq 1024$ contributors \cite{bowegabizonmiers}. Knowledge
soundness requires the BLS12-381 discrete log assumption and is
therefore \emph{not} post-quantum: a quantum adversary running Shor
on $\mathbb{F}_{q^{12}}$ can extract trapdoors and forge
$\pi_{\mathrm{ZK}}$. However, $\pi_{\mathrm{ZK}}$ commits to the
public-input hash $h_x$ which in turn commits to the validator set
and round digest; forging $\pi_{\mathrm{ZK}}$ does not produce
forged Pulsar / Corona / Magnetar signatures, so the conjunctive
composition still requires breaking those PQ legs to forge the
cert as a whole (Theorem~\ref{thm:composition}).

\subsection{Hash functions and domain separation}

All hash invocations use domain-separated contexts per FIPS 204 \S5
and SP 800-185 (cSHAKE / KMAC). Table~\ref{tab:domsep} lists the
context strings used by the QuasarCert legs.

\begin{table}[ht]
\centering
\small
\begin{tabular}{@{}lll@{}}
\toprule
\textbf{Context string} & \textbf{Hash} & \textbf{Used by} \\
\midrule
\texttt{lux-quasar-cert-v1} & SHAKE-256 & Pulsar leg of QuasarCert \\
\texttt{lux-corona-cert-v1} & SHAKE-256 & Corona leg of QuasarCert \\
\texttt{lux-magnetar-cert-v1} & SHAKE-256 & Magnetar leg of QuasarCert \\
\texttt{lux-quasar-zchain-witness-v1} & SHAKE-256 & per-validator ML-DSA into Groth16 \\
\texttt{lux-quasar-bls-domain-v1} & SHA-256 & BLS hash-to-curve $H_{\mathbb{G}_1}$ \\
\texttt{lux-quasar-round-digest-v1} & SHA3-256 & round-digest $\mu$ construction \\
\texttt{lux-quasar-validator-root-v1} & SHA3-256 & validator-set Merkle root \\
\bottomrule
\end{tabular}
\caption{Domain-separation contexts used by QuasarCert legs.
Mismatches across legs prevent cross-leg signature transfer (e.g.,
a BLS signature cannot be replayed as a Pulsar input even if the
same byte string is presented).}
\label{tab:domsep}
\end{table}
